% GENERATED FILE. Edit canonical lessons, then run npm run generate:books.
% canonical-source-hash: fnv1a64:f4283087b386cfaf
% canonical-lessons: BN-C221-besorkari, BN-C221-jatiyo, BN-C221-antorjatik, BN-C221-shohure, BN-C221-niyomito

\chapter{Private, National, International, Urban, Regular}
\label{ch:bn-private-national-international-urban-regular}
\hlchaptermodality{221}

\begin{chapteropening}
\textbf{By the end of this chapter:} \emph{I can say private, national, international, urban and regular.}

Complete the last lesson of chapter 221: I can say private, national, international, urban and regular.
\end{chapteropening}

\section[\texorpdfstring{\bn{বেসরকারি} (besôrkāri)}{বেসরকারি (besôrkāri)}]{\bn{বেসরকারি} (besôrkāri) — private}

\label{lesson:BN-C221-besorkari}



\begin{quote}
Before the new one: say the Bengali for excellent, then the Bengali for previous.
\end{quote}

\subsection*{You'll want to know: \bn{বেসরকারি}}
\textbf{\bn{বেসরকারি}} — \emph{besôrkāri} — \textquotedblleft{}private\textquotedblright{}.

\begin{grammarlens}[title={What you've built}]
One more word for this chapter.
\end{grammarlens}

\subsection*{Your turn}
\begin{itemize}
\raggedright
  \item \emph{Say it:} \emph{besôrkāri}
  \item \emph{Say it:} \emph{besôrkāri}, once more
  \item \emph{Say it:} say \emph{besôrkāri}
  \item \emph{Recall:} say the Bengali for excellent, then the Bengali for previous, then say \emph{besôrkāri} again
\end{itemize}

\begin{tcolorbox}[breakable,colback=teal!4,colframe=teal!35!black,title={Before you move on}]
What does \bn{বেসরকারি} mean? (\textquotedblleft{}Private\textquotedblright{}.) Say it once more.
\end{tcolorbox}
\section[\texorpdfstring{\bn{জাতীয়} (jātīyô)}{জাতীয় (jātīyô)}]{\bn{জাতীয়} (jātīyô) — national}

\label{lesson:BN-C221-jatiyo}



\begin{quote}
Before the new one: say the Bengali for previous, then the Bengali for private.
\end{quote}

\subsection*{You'll want to know: \bn{জাতীয়}}
\textbf{\bn{জাতীয়}} — \emph{jātīyô} — \textquotedblleft{}national\textquotedblright{}.

\begin{grammarlens}[title={What you've built}]
One more word for this chapter.
\end{grammarlens}

\subsection*{Your turn}
\begin{itemize}
\raggedright
  \item \emph{Say it:} \emph{jātīyô}
  \item \emph{Say it:} \emph{jātīyô}, once more
  \item \emph{Say it:} say \emph{jātīyô}
  \item \emph{Recall:} say the Bengali for previous, then the Bengali for private, then say \emph{jātīyô} again
\end{itemize}

\begin{tcolorbox}[breakable,colback=teal!4,colframe=teal!35!black,title={Before you move on}]
What does \bn{জাতীয়} mean? (\textquotedblleft{}National\textquotedblright{}.) Say it once more.
\end{tcolorbox}
\section[\texorpdfstring{\bn{আন্তর্জাতিক} (āntôrjātik)}{আন্তর্জাতিক (āntôrjātik)}]{\bn{আন্তর্জাতিক} (āntôrjātik) — international}

\label{lesson:BN-C221-antorjatik}



\begin{quote}
Before the new one: say the Bengali for private, then the Bengali for national.
\end{quote}

\subsection*{You'll want to know: \bn{আন্তর্জাতিক}}
\textbf{\bn{আন্তর্জাতিক}} — \emph{āntôrjātik} — \textquotedblleft{}international\textquotedblright{}.

\begin{grammarlens}[title={What you've built}]
One more word for this chapter.
\end{grammarlens}

\subsection*{Your turn}
\begin{itemize}
\raggedright
  \item \emph{Say it:} \emph{āntôrjātik}
  \item \emph{Say it:} \emph{āntôrjātik}, once more
  \item \emph{Say it:} say \emph{āntôrjātik}
  \item \emph{Recall:} say the Bengali for private, then the Bengali for national, then say \emph{āntôrjātik} again
\end{itemize}

\begin{tcolorbox}[breakable,colback=teal!4,colframe=teal!35!black,title={Before you move on}]
What does \bn{আন্তর্জাতিক} mean? (\textquotedblleft{}International\textquotedblright{}.) Say it once more.
\end{tcolorbox}
\section[\texorpdfstring{\bn{শহুরে} (shôhure)}{শহুরে (shôhure)}]{\bn{শহুরে} (shôhure) — urban}

\label{lesson:BN-C221-shohure}



\begin{quote}
Before the new one: say the Bengali for national, then the Bengali for international.
\end{quote}

\subsection*{You'll want to know: \bn{শহুরে}}
\textbf{\bn{শহুরে}} — \emph{shôhure} — \textquotedblleft{}urban\textquotedblright{}.

\begin{grammarlens}[title={What you've built}]
One more word for this chapter.
\end{grammarlens}

\subsection*{Your turn}
\begin{itemize}
\raggedright
  \item \emph{Say it:} \emph{shôhure}
  \item \emph{Say it:} \emph{shôhure}, once more
  \item \emph{Say it:} say \emph{shôhure}
  \item \emph{Recall:} say the Bengali for national, then the Bengali for international, then say \emph{shôhure} again
\end{itemize}

\begin{tcolorbox}[breakable,colback=teal!4,colframe=teal!35!black,title={Before you move on}]
What does \bn{শহুরে} mean? (\textquotedblleft{}Urban\textquotedblright{}.) Say it once more.
\end{tcolorbox}
\section[\texorpdfstring{\bn{নিয়মিত} (niyômitô)}{নিয়মিত (niyômitô)}]{\bn{নিয়মিত} (niyômitô) — regular}

\label{lesson:BN-C221-niyomito}



\begin{quote}
Before the new one: say the Bengali for international, then the Bengali for urban.
\end{quote}

\subsection*{You'll want to know: \bn{নিয়মিত}}
\textbf{\bn{নিয়মিত}} — \emph{niyômitô} — \textquotedblleft{}regular\textquotedblright{}.

\begin{grammarlens}[title={What you've built}]
One more word for this chapter.
\end{grammarlens}

\subsection*{Your turn}
\begin{itemize}
\raggedright
  \item \emph{Say it:} \emph{niyômitô}
  \item \emph{Say it:} \emph{niyômitô}, once more
  \item \emph{Say it:} say \emph{niyômitô}
  \item \emph{Recall:} say the Bengali for international, then the Bengali for urban, then say \emph{niyômitô} again
\end{itemize}

\begin{tcolorbox}[breakable,colback=teal!4,colframe=teal!35!black,title={Before you move on}]
What does \bn{নিয়মিত} mean? (\textquotedblleft{}Regular\textquotedblright{}.) Say it once more.
\end{tcolorbox}
